\section{Moving contacts and the normalization of the failure divisor}
\label{sec:moving-contact-normalization}
Let $X_d^\circ=X_{d,2}^\circ$ be the relative generating Grassmannian
in Theorem~\ref{thm:sharp-conductor}. It is smooth and irreducible
of dimension $3d-4$. The fixed-divisor schemes of the preceding
section are its scheme-theoretic fibres, not its irreducible
components. For $m\ge d-1$ denote its failure scheme by $Y_m$.

\subsection{Components over every multiplicity stratum}
For a partition $\boldsymbol d=(d_1,\ldots,d_s)$ of $d$, let
$B_{\boldsymbol d}\subset\operatorname{Sym}^d\Pj^1$ be the locally
closed stratum of divisors with exactly this multiplicity partition,
and restrict $X_d^\circ$ to it. Let $M$ be the set of distinct
multiplicities, and let $r_e=\#\{i:d_i=e\}$ for $e\in M$.

\begin{corollary}[Complete component descent on multiplicity strata]
\label{cor:contact-multiplicity-strata}
For $m\ge d-1$, the failure scheme over $B_{\boldsymbol d}$ has
\begin{equation}\label{eq:contact-stratum-component-count}
 \#\{e\in M:e\ge2\}+\binom{\#M}{2}
                         +\#\{e\in M:r_e\ge2\}
\end{equation}
irreducible components, all of dimension $s+2d-5$, and no embedded
associated points. A ramification component is indexed by each
$e\ge2$ in $M$ and has primary multiplicity $\binom e2$. An
identification component is indexed by each unordered pair of
multiplicities $e,f$ that occurs at distinct support points and has
primary multiplicity $ef$. The ideal of the scheme is the
intersection of these prime-divisor ideals raised to the stated
powers. On the labelled etale cover its full local equations and
nilradical layers are exactly those of
Theorem~\ref{thm:contact-layers}.
\end{corollary}
\begin{proof}
The ordered configuration space of $s$ distinct points of $\Pj^1$
is an integral finite etale cover of $B_{\boldsymbol d}$, with
deck group $\prod_{e\in M}\mathfrak S_{r_e}$. The relative
first-jet and two-point restriction maps are surjective over this
cover. Consequently each labelled $T_i$ and $E_{ij}$ is an
irreducible relative rank-one Schubert divisor. The local
factorization \eqref{eq:contact-confluent-determinant} holds with
the support coordinates free, so there are no further components.
The deck group permutes $T_i$ exactly according to the value $d_i$,
and permutes $E_{ij}$ exactly according to the unordered pair
$\{d_i,d_j\}$. The images of the orbit unions are precisely the
irreducible components downstairs: each is the image of an integral
labelled component, and two such images agree exactly for components
in the same orbit. This gives \eqref{eq:contact-stratum-component-count}.

Lemma~\ref{lem:etale-primary-descent} applies to this regular ambient
space and its invariant weighted divisors, including all branch
intersections. It proves that each component downstairs is a prime
Cartier divisor; irreducibility here does not assert smoothness or
geometric unibranchedness. Pulling its reduced ideal back to the labelled cover
gives the reduced product of its labelled branch equations. Their
multiplicity is constant on an orbit, so the primary ideal identity
in \eqref{eq:contact-full-primary} descends faithfully flatly. This
also proves absence of embedded points and gives the stated local
layers. The base stratum has dimension $s$, giving the dimension
formula. A descended component can have intersecting local branches;
smoothness of each labelled divisor is not a claim that every
descended component is smooth.
\end{proof}

\begin{theorem}[An integral failure divisor with smooth normalization]
\label{thm:moving-normalization}
For $d\ge4$, $n\ge2d-2$, and $m\ge d-1$, the schemes $Y_m$
are all the same integral Cartier divisor $Y$ of dimension $3d-5$.
In particular $Y$ has no embedded associated points. Its normalization is
\begin{equation}\label{eq:contact-normalization-incidence}
 \widetilde Y=
 \{(D,A,Z): (D,A)\in X_d^\circ,\ Z\le D,\ \deg Z=2,
                  \rank(A\longrightarrow H^0(\mathcal O_Z(n)))=1\}.
\end{equation}
The forgetful map $\widetilde Y\to Y$ is finite and birational,
and $\widetilde Y$ is smooth and irreducible.

More explicitly, on an affine support chart write $D=(g(t))$ with
$g$ monic of degree $d$, and on a frame chart put $v=b/a\in\C[t]/(g)$,
represented by a polynomial of degree less than $d$. The failure
equation is
\begin{equation}\label{eq:moving-power-determinant}
 \Delta(g,v)=\det[1,v,\ldots,v^{d-1}]_{\C[t]/(g)}=0.
\end{equation}
The normalization on this chart is parametrized by
\begin{equation}\label{eq:moving-explicit-normalization}
 g=hq,\qquad v=\lambda+hw,
\end{equation}
where $h,q$ are monic of degrees $2,d-2$, respectively,
$\lambda\in\C$, and $\deg w\le d-3$. The scalar-pencil locus
is omitted; the unit $a$ is a free smooth parameter.
\end{theorem}
\begin{proof}
Theorem~\ref{thm:sharp-conductor} applies on the whole relative
space. The proof of Lemma~\ref{lem:contact-power-basis} applies over
the polynomial coefficient ring of a varying monic $g$, not only
on its fibres. Thus every $Y_m$ has the identical ideal
\eqref{eq:moving-power-determinant} on a smooth covering by frame
charts. This determinant is not identically zero: $v=t$ makes its
matrix the identity for every $g$.

The space of pairs $Z\le D$ with lengths $2,d$ is
$\operatorname{Sym}^2\Pj^1\times\operatorname{Sym}^{d-2}\Pj^1$;
the map sends $(Z,Q)$ to $D=Z+Q$. It is finite. Indeed it is proper,
and a fixed effective divisor has only finitely many degree-two
subdivisors, obtained by choosing multiplicities at its finitely
many support points. Restriction $E_{n,D}\to E_{n,Z}$ is a
surjection of vector bundles. The rank-at-most-one condition on a
pencil is a relative Schubert divisor; its rank-zero part is
excluded by the generating hypothesis. This makes
\eqref{eq:contact-normalization-incidence} a closed subscheme of
the finite base change to $X_d^\circ$, so its projection is finite.

Smoothness and irreducibility can be checked explicitly. A divisor
with finite support is contained in an affine coordinate chart. On
that chart $Z$ is defined by a unique monic quadratic $h$, and
$Z\le D$ means $g=hq$. Since $a$ is a unit, rank one on $Z$ is
equivalent to $v\bmod h$ being a scalar $\lambda$. Division by the
monic $h$ gives uniquely $v=\lambda+hw$, with the stated degree
bound. Thus the framed incidence is an open part of affine space
with coordinates the coefficients of $h,q,w$, the scalar $\lambda$,
and the unit $a$. This proves smoothness. Globally irreducibility
also follows from the relative rank-one Schubert description over
the integral pair-divisor base; each fibre is irreducible, and its
generating open is nonempty. These charts describe the same
incidence, so no extra component appears at infinity.

An element $v$ fails to generate the length-$d$ algebra
$R=H^0(\mathcal O_D)$ precisely when it either takes the same value
at two distinct support points or has zero first derivative at a
multiple support point. This follows at once from the minimal
polynomial \eqref{eq:contact-minimal-polynomial}: its degree is $d$
if and only if the support values are distinct and each local
$\ell_i=d_i$. These two failure mechanisms are precisely the
existence of a length-two $Z\le D$ on which $v$ is constant.
Therefore the image of the finite incidence is the whole support
of \eqref{eq:moving-power-determinant}. That support is irreducible.

To determine its multiplicity, work over divisors with distinct
support, label their roots after an etale base change, and use
evaluation coordinates for $v$. The fixed-fibre formula becomes a
product of the distinct differences $\lambda_j-\lambda_i$. At a
point with exactly one equal pair the determinant has a nonzero
linear derivative. Hence its multiplicity along its unique
irreducible support is one. Since the ambient space is smooth, a
Cartier hypersurface has no embedded associated points; a
generically reduced such hypersurface is reduced. One may see the
last assertions locally from unique factorization in the regular
local ring, or from the Cohen--Macaulay and $R_0+S_1$ criterion
\cite[Tag 031R]{Stacks}. Thus $Y$ is integral.

Over the dense open subset with reduced $D$ and precisely one equal
pair there is exactly one possible $Z$. The finite projection is
therefore birational in characteristic zero. Its source is smooth
and hence normal, so it is the normalization. The affine
parametrization is exactly \eqref{eq:moving-explicit-normalization},
as already proved. The dimension of the divisor is $3d-5$.
\end{proof}

\begin{proposition}[Exact singular-support test across collisions]
\label{prop:moving-singular-test}
At a closed point $(D,A)\in Y$, call a length-two $Z\le D$
\emph{scalar} if $A$ has rank one on $Z$. The point is smooth on
$Y$ if and only if there is precisely one scalar $Z$ and the
following linear system has no nonzero solution. In the associated
factorization \eqref{eq:moving-explicit-normalization}, the unknown
is $r(t)$ of degree less than two, and the system is
\begin{equation}\label{eq:moving-singular-system}
                 h\mid qr,\qquad wr\bmod h\text{ is constant}.
\end{equation}
In particular, this criterion is effective at nonreduced $D$ and
is not merely the distinct-support pair-count criterion. The full
singular scheme on a frame chart has ideal
\begin{equation}\label{eq:moving-jacobian-ideal}
 (\Delta,\partial\Delta/\partial g_0,\ldots,
          \partial\Delta/\partial g_{d-1},
          \partial\Delta/\partial v_0,\ldots,
          \partial\Delta/\partial v_{d-1}).
\end{equation}
\end{proposition}
\begin{proof}
If the normalization has more than one point above $(D,A)$, the
local ring of $Y$ cannot be regular. Suppose there is just one.
We first justify the differential criterion for a finite
normalization $\nu:N\to Y$ with $N$ smooth. A regular local ring
is normal, so a smooth point of $Y$ has an isomorphic normalization
near it. Conversely let $A_0\subset B_0$ be the local finite
extension at a point with a unique preimage and common residue
field $\C$. Injectivity of $d\nu$ means
\[
 \mathfrak m_{B_0}/(\mathfrak m_{B_0}^2+
                         \mathfrak m_{A_0}B_0)=0.
\]
Nakayama's lemma first gives
$\mathfrak m_{B_0}=\mathfrak m_{A_0}B_0$. It then gives, for the
finite $A_0$-module $B_0$, that $1$ generates $B_0$, since the
quotient by $\mathfrak m_{A_0}$ is $\C$. The injective map
$A_0\to B_0$ is consequently surjective. Thus $A_0=B_0$ is regular.
Local finiteness here follows by localizing the finite morphism;
uniqueness of the preimage makes the resulting semilocal ring local.

In the affine normalization coordinates, a tangent vector with
$\dot g=\dot v=0$ satisfies
\[
 0=q\dot h+h\dot q,\qquad
 0=\dot\lambda+w\dot h+h\dot w.
\]
Set $r=\dot h$, so $\deg r<2$ because $h$ is monic. The first
equation is soluble exactly when $h\mid qr$, and then determines
$\dot q=-qr/h$. The second is soluble exactly when the remainder
of $wr$ modulo $h$ is constant; it determines $\dot\lambda$ as
the negative of that constant and determines $\dot w$ by division.
The resulting degrees satisfy the required bounds. If $r=0$ all
four variations are zero. Thus the differential has nonzero kernel
exactly when \eqref{eq:moving-singular-system} has a nonzero solution.
The unit frame parameter has zero variation in a kernel and creates
no additional one. Smooth frame descent proves the criterion on
$Y$. Finally \eqref{eq:moving-jacobian-ideal} is the hypersurface
Jacobian presentation in its smooth affine ambient space.
\end{proof}


\begin{remark}[A smooth total point with a triple primary fibre]
\label{rem:contact-smooth-total-triple-fibre}
Take $d=4$, $g=t^3(t-1)$, and $A=\langle1,t^2\rangle$ on
$D=(g)$. The only scalar length-two subdivisor is $h=t^2$.
In \eqref{eq:moving-explicit-normalization} one has
$q=t(t-1)$ and $w=1$. For $r=r_0+r_1t$, divisibility
$h\mid qr$ forces $r_0=0$, while constancy of $wr\bmod h$
forces $r_1=0$. Hence the total divisor is smooth here. On the
fixed-divisor fibre, the two support values are distinct and
\eqref{eq:contact-confluent-determinant} is a unit times $c_{11}^3$.
Its exact local nilpotency index is three. Thus even a smooth point
of the integral total divisor can have a nonreduced primary fibre.
\end{remark}

\begin{proposition}[Discriminants and specialization]
\label{prop:contact-discriminant}
Let $\chi_v(T)$ be the characteristic polynomial of multiplication
by $v$ in $\C[t]/(g)$. With the usual monic polynomial
normalization of discriminants,
\begin{equation}\label{eq:contact-discriminant-identity}
 \operatorname{disc}(\chi_v)=
                     \operatorname{disc}(g)\,\Delta(g,v)^2.
\end{equation}
This is a polynomial identity on the entire coefficient space.
The equation $\Delta=0$, not $\operatorname{disc}(\chi_v)=0$,
defines the failure divisor, including when $g$ has multiple roots.
Its fibre over each fixed $D$ is exactly the primary scheme in
\eqref{eq:contact-full-primary}.
\end{proposition}
\begin{proof}
When $g$ has distinct roots $t_1,\ldots,t_d$, the evaluation
matrix changes the basis $1,t,\ldots,t^{d-1}$ to point values.
Therefore
\[
 \Delta(g,v)=\frac{\prod_{i<j}(v(t_j)-v(t_i))}
                   {\prod_{i<j}(t_j-t_i)}.
\]
Squaring gives \eqref{eq:contact-discriminant-identity}, because
$\chi_v(T)=\prod_i(T-v(t_i))$ on this open set. Both sides are
polynomial in the coefficients, so the identity holds everywhere.
For $v=t$ the determinant is one even when $g$ is nonreduced; this
also shows why the factor $\operatorname{disc}(g)$ must not be
included in the failure equation. The scheme-theoretic fibre
statement follows from Fitting base change in
Theorem~\ref{thm:sharp-conductor}, or by specializing the
polynomial determinant and applying
\eqref{eq:contact-confluent-determinant}.
\end{proof}

\begin{proof}[Proof of Theorem~\ref{thm:contact-main}]
The fixed-contact assertions are
Theorems~\ref{thm:contact-factorization} and~\ref{thm:contact-layers}.
The moving-contact assertions are
Theorem~\ref{thm:moving-normalization} and
Proposition~\ref{prop:moving-singular-test}. All four concern the
original multiplication map, by the coherent cokernel isomorphism
of Theorem~\ref{thm:sharp-conductor}.
\end{proof}
